\section{Gaussian enlargement and the Rayleigh comparison}
\label{S:actual-gaussian-tubes}

We first state the Gaussian geometry with its exact covariance scale. Let $H$ be
real separable Hilbert, $W$ its given isonormal process, and $X$ a strongly
measurable centered Gaussian variable in a real separable Banach space
$\mathscr{B}$, with $\E\|X\|^2 < \infty$. Assume every continuous linear
coordinate has the form $\ell(X) = W(k_\ell)$, $k_\ell \in H$. Fix $0 < \delta <
1$ and write
\begin{equation}\label{E:tube-covariance-scale}
  Ch=\E[XW(h)]=\sqrt\delta\,\mathcal{M}_\delta h,\qquad
  H_{\rm seen}=\overline{\operatorname{span}\{k_\ell\}},\qquad
  E_0=\overline{C(H)} .
\end{equation}
Assume $\mathcal{M}_\delta:H \to \mathscr{B}$ is compact. All Hilbert balls
below are closed.

\subsection{The Wiener-space realization}

\begin{proposition}[First Wiener chaos realization on the support]
  \label{P:tube-wiener-realization}
  The variable $X$ lies in $E_0$ almost surely, with a centered Radon Gaussian
  law of full support there. Its Cameron space and norm are
  \begin{equation}\label{E:tube-cameron-space}
    H_{\rm CM}=C(H_{\rm seen}),\qquad
    \|Ch\|_{\rm CM}=\|P_{\rm seen}h\|_H .
  \end{equation}
  This support triple is an abstract Wiener realization, including the
  measurable-norm condition.
\end{proposition}

\begin{proof}
  The zero law has $C = 0$ and $E_0 = \{0\}$ and is immediate. A countable
  norming family exists in $\mathscr{B}^*$: choose a dense sequence on the unit
  sphere and apply Hahn--Banach at each point. Approximation gives $\|x\| =
  \sup_j|\ell_j(x)|$. Norms of differences from a countable dense set show that
  these coordinates generate the Borel sets. Consequently, $X$ is measurable in
  the completion of $\sigma(W(H))$. Fix an orthonormal basis $(e_j)$ of $H$ and
  set $\mathcal{G}_j = \sigma(W(e_1),\ldots,W(e_j))$. In finite dimension these
  sums and projections stabilize at the dimension. Bochner conditional
  expectations need no Radon--Nikodym property here: define them first on
  finite-valued simple variables using scalar conditional expectations, then
  extend by the $L^2$ norm contraction. Simple approximation supplies their
  scalar testing identities. Gaussian regression and the countable norming
  family give
  \begin{equation}\label{E:tube-finite-regression}
    X_j=\E[X\mid\mathcal{G}_j]
    =\sum_{l\le j}W(e_l)Ce_l,\qquad X_j\longrightarrow X
    \quad\hbox{in }L^2(\mathscr{B}).
  \end{equation}
  To justify convergence, finite-coordinate scalar events approximate all events
  in the generated sigma field: the class of events whose indicators are
  approximable in $L^2$ is a monotone class containing the finite-coordinate
  algebra. Approximate a Bochner variable by simple variables and use the
  conditional $L^2$ contraction.

  Bounded testing gives $\ell(Ch) = \langle k_\ell,h\rangle$. Thus, $\ker C =
  H_{\rm seen}^\perp$, and $C$ is injective on $H_{\rm seen}$. The closed first
  Wiener chaos generated by $X$ is exactly $W(H_{\rm seen})$. Its covariance
  image therefore has the first Wiener chaos norm
  in~\eqref{E:tube-cameron-space}; unused directions are removed by the
  orthogonal projection, without an inverse covariance. This is the Cameron
  image under the first Wiener chaos definition. The finite sums
  in~\eqref{E:tube-finite-regression} lie in $E_0$, so $L^2$ convergence and an
  almost-sure subsequence put $X$ in that closed space almost surely.
  Functionals on $E_0$ extend to $\mathscr{B}$ by Hahn--Banach, so restriction
  preserves the scalar laws and first chaos.

  Its law is tight. For each radius $2^{-j}$ choose a finite union of balls of
  mass greater than $1-\varepsilon2^{-j-1}$, using a countable cover of the
  separable space. The intersection of the corresponding finite closed unions is
  closed and totally bounded and has mass at least $1-\varepsilon$. Completeness
  makes it compact. Metric Borel regularity on these compact sets, followed by
  tightness, gives a Radon law. To prove full support, the tail $X-X_j$ is
  independent of the finite Gaussian sum: every tail coordinate has label
  $(I-P_j)k_\ell$, orthogonal to $e_1,\ldots,e_j$. Joint Gaussianity and
  countable Borel coordinates give Banach tail independence. Given $h$ and
  $\varepsilon > 0$, choose $j$ with $\|CP_jh-Ch\| < \varepsilon/3$ and
  $\mathbb{P}(\|X-X_j\| < \varepsilon/3) > 0$. The finite sum has positive
  probability in every ball about $CP_jh$, even if its covariance is singular:
  the preimage under its finite linear map contains an open ball about the
  coefficient vector of $P_jh$. Independence proves positivity of every ball
  about $Ch$. Closure proves that the support is all of $E_0$.

  We check the measurable norm explicitly. In $H_{\rm seen}$ put $V_j =
  \operatorname{span}\{P_{\rm seen}e_l:l\le j\}$. Any finite orthonormal family
  $(h_k)$ in $V_j^\perp$ is also orthogonal to $e_1,\ldots,e_j$. The same
  regression and Bochner Jensen argument gives
  \[
    \E[X-X_j\mid W(h_k),\hbox{ all }k]
    =\sum_k W(h_k)Ch_k,\qquad
    \E\left\|\sum_kW(h_k)Ch_k\right\|^2
    \le\E\|X-X_j\|^2 .
  \]
  For any $\varepsilon > 0$ choose $j$ so that the right side is less than
  $\varepsilon^3$. Markov bounds the probability that this finite sum has norm
  exceeding $\varepsilon$ by $\varepsilon$, uniformly over every such finite
  orthonormal family. Since $C:H_{\rm seen} \to H_{\rm CM}$ is isometric, the
  excluded finite-dimensional Cameron subspace is $C(V_j)$. This is the
  measurable-norm tail condition for the dense Hilbert image, whose Banach
  completion is $E_0$. The characteristic functional of the law is $\E
  e^{i\ell(X)} = \exp(-\|k_\ell\|_H^2/2)$. These facts give the stated abstract
  Wiener realization. The zero space has the direct zero-law interpretation.
  This proves Proposition~\ref{P:tube-wiener-realization}.
\end{proof}

\subsection{Closed Gaussian enlargements and the quadratic cost}
\label{S:gaussian-isoperimetry-input}

The Gaussian isoperimetric theorem is due to Borell and, independently,
Sudakov--Tsirelson~\cites{borell:75:brunn-minkowski, sudakov.cirel-son:74:extremal}.
We use its set-enlargement form from
Friz and Hairer~\cite{friz.hairer:20:course}*{Theorem~11.6, p.~190}, with footnote~3 on
p.~190; p.~191 supplies the following context. For an abstract Wiener space, a
Borel $A$ of positive mass and its Cameron unit ball $\mathcal{K}$, it states
\[
  \mu((A+t\mathcal{K})^c)
  \le\overline{\Phi}_{\rm N}(\widehat{a}+t),\qquad
  \widehat{a}=\Phi_{\rm N}^{-1}(\mu(A)),\quad t\ge0 .
\]
Here $\Phi_{\rm N}$ is the standard normal distribution function; the same page
gives $\overline{\Phi}_{\rm N}(y) \le e^{-y^2/2}$ for $y \ge 0$. This theorem is
imported, not proved here. It concerns Gaussian input sets, not Gaussianity of a
nonlinear wave. We verify Borel measurability of every enlargement we use.

\begin{lemma}[Compact images and attained tube radii]
  \label{L:tube-closed-geometry}
  Suppose the law of $X$ is nonzero and put
  \[
    b_\delta=2\E\|X\|,\qquad
    A_\delta=\{x\in E_0:\|x\|\le b_\delta\},\qquad
    \mathcal{T}_r=A_\delta+\mathcal{M}_\delta(\overline{B}_H(r)).
  \]
  Every $\mathcal{T}_r$ is closed. The extended distance $d_\delta(x) =
  \inf\{r\ge0:x\in\mathcal{T}_r\}$ satisfies
  \begin{equation}\label{E:tube-attained-radius}
    \{d_\delta\le r\}=\mathcal{T}_r,\qquad
    \mathbb{P}(d_\delta(X)>r)\le e^{-r^2/(2\delta)},\quad r\ge0 .
  \end{equation}
  In particular, $d_\delta$ is lower semicontinuous and finite almost surely.
\end{lemma}

\begin{proof}
  A compact map from a separable Hilbert space sends a closed bounded ball to a
  compact set, including closedness of the image. Indeed, diagonal orthonormal
  coordinates and Fatou give a weakly convergent subsequence of any bounded
  sequence, with limit still in the ball. Compactness gives norm-convergent
  image subsequences, and bounded scalar testing identifies every limit with the
  image of that weak limit. This proves sequential compactness. Closed
  $A_\delta$ plus a compact image is closed: along a converging sequence of
  sums, extract the compact component and pass the residual into the closed base
  set.

  Membership in $\mathcal{T}_r$ implies $d_\delta \le r$. Conversely, if
  $d_\delta(x) \le r$, choose $h_j$ with $\|h_j\| \le r+1/j$ and
  $x-\mathcal{M}_\delta h_j \in A_\delta$. A weak subsequence has limit $h$ of
  norm at most $r$; compactness gives $\mathcal{M}_\delta h_j \to
  \mathcal{M}_\delta h$. Closedness of the residual proves membership in
  $\mathcal{T}_r$. This also covers $r = 0$ and proves least finite-radius
  attainment. Closed sublevels give lower semicontinuity and Borel
  measurability.

  The nonzero law has $b_\delta > 0$, and Markov gives $\mu(A_\delta) \ge 1/2$.
  Projection onto $H_{\rm seen}$ removes all unused directions in the original
  covariance space, so~\eqref{E:tube-cameron-space} gives the exact set
  identities
  \[
    \mathcal{K}=C(\overline{B}_H(1))
    =\sqrt\delta\,\mathcal{M}_\delta(\overline{B}_H(1)),\qquad
    \mathcal{T}_r=A_\delta+(r/\sqrt\delta)\mathcal{K} .
  \]
  These are Borel enlargements by the preceding closedness. The imported theorem
  applies to the realization in Proposition~\ref{P:tube-wiener-realization}.
  Since $\widehat{a} \ge 0$, it yields~\eqref{E:tube-attained-radius}. The value
  $\widehat{a} = +\infty$ is allowed when $\mu(A_\delta) = 1$. The tail tends to
  zero, proving almost-sure finiteness. The Cameron embedding does not change
  the factor one half in the Gaussian cost. This proves
  Lemma~\ref{L:tube-closed-geometry}.
\end{proof}

\begin{figure}
  \centering
  % Schematic geometry only; the radius belongs to the control space.
\begin{tikzpicture}[x=1cm,y=1cm,font=\normalsize]
  \filldraw[draw=hamteal,fill=hamteal!5,thick,dashed]
  (2.7,2) ellipse (2.3 and 1.4);
  \filldraw[draw=hamblue,fill=hamblue!12,thick]
  (2.1,1.8) ellipse (.85 and .55);
  \node at (2.1,1.8) {$A_\delta$};
  \node[anchor=south] at (2.7,3.65)
  {$A_\delta+\mathcal{M}_\delta(\overline{B}_{H_I}(r))$};
  \draw[hamteal] (2.7,3.6)--(2.7,3.4);
  \fill (6.3,2.5) circle (1.7pt);
  \node[align=center,anchor=north] at (6.3,2.25)
  {outside the\\enlarged set};
  \node[ham/box,font=\normalsize,text width=4.6cm] at (10.05,2.6)
  {Probability of the complement\\
  $\le \exp\{-r^2/(2\delta)\}$};
  \node[align=center,text width=5cm] at (10.05,.85)
  {$r$ is a control-space radius.\\
  It is not Euclidean distance\\in the picture.};
\end{tikzpicture}

  \caption[Gaussian enlargement by an image of a Hilbert ball]{A schematic
    finite-dimensional projection of enlargement for the nonzero Gaussian law in
    the cited lemma. The map $\mathcal{M}_\delta$ is compact, so the image of a
    closed control ball is compact; the set $A_\delta$ itself need not be
    compact. No ambient Lebesgue density or Euclidean distance formula is being
    assumed. The cost uses one global control norm. See
    \ProofRef{Lemma}{L:tube-closed-geometry} and
    \ProofRef{Proposition}{P:tube-rayleigh-transfer}.}
  \label{Fig:gaussian-enlargement}
\end{figure}


\begin{proposition}[Comparison with the Rayleigh integral, including jumps and zero law]
  \label{P:tube-rayleigh-transfer}
  Let $F:E_0 \to \R$ be Borel and let $G:[0,\infty) \to \R$ be Borel and
  nondecreasing. If $F(x) \le G(r)$ for every $x \in \mathcal{T}_r$ and every $r
  \ge 0$, then
  \begin{equation}\label{E:tube-rayleigh-transfer}
    \E e^{F(X)/\delta}
    \le\int_0^\infty\frac r\delta
    \exp\left(\frac{G(r)-r^2/2}{\delta}\right)\ud r .
  \end{equation}
  For zero law use $E_0 = \{0\}$, $C = \mathcal{M}_\delta = 0$; the same
  conclusion holds when $F(0) \le G(0)$.
\end{proposition}

\begin{proof}
  The auxiliary Rayleigh variable $\mathcal{R}_\delta$ with density
  $(r/\delta)e^{-r^2/(2\delta)}$, $r > 0$, has mass one and survival function
  $e^{-r^2/(2\delta)}$ by direct integration. It is a scalar stochastic
  dominator, not a norm of a Gaussian variable in the covariance Hilbert space.
  For every level $t \ge 0$, the upper level set of $e^{G(r)/\delta}$ is an
  upper interval, possibly empty or the whole halfline. An open endpoint is
  controlled by~\eqref{E:tube-attained-radius}. At a positive closed endpoint
  $c$, let $s\uparrow c$:
  \[
    \mathbb{P}(d_\delta(X)\ge c)
    =\lim_{s\uparrow c}\mathbb{P}(d_\delta(X)>s)
    \le e^{-c^2/(2\delta)}
    =\mathbb{P}(\mathcal{R}_\delta\ge c).
  \]
  Rayleigh has no atoms, and at the closed endpoint zero both probability bounds
  are one. On the null event $\{d_\delta(X)=\infty\}$ define
  $e^{G(d_\delta(X))/\delta}$ to be one. No value of $G$ at infinity is needed
  and no expectation changes. Layer cake and Tonelli therefore order the
  expectations of $e^{G(d_\delta(X))/\delta}$ and
  $e^{G(\mathcal{R}_\delta)/\delta}$, even when $G$ jumps or expectations are
  infinite. Least-radius attainment gives $F(X) \le G(d_\delta(X))$ almost
  surely. No measurable selection of a control or residual is needed. This
  proves~\eqref{E:tube-rayleigh-transfer}. For zero law, $X = 0$ and $F(0) \le
  G(0) \le G(r)$; normalization of the Rayleigh density proves the claim
  directly. In particular, $F(0)$ need not be zero. This proves
  Proposition~\ref{P:tube-rayleigh-transfer}.
\end{proof}

\subsection{The wave tuple and its globally defined endpoint}

Retain all original covariances and amplitudes. Use the exponents and fixed even
$Q$ in~\eqref{E:radius-anchor-parameters}--\eqref{E:radius-anchor-Q}; put $\nu =
1/\theta$ and choose $(1-\theta)^{-1} < r_t < 2/a$, $\nu \le p < 2 < r_t$. All
these intervals are nonempty by the strict subcritical assumptions. Fix integer
$B \ge 1$ and the true partition $\delta,N,s_m,\ell_m$ from
Section~\ref{S:actual-cameron-mean}. Set $L_\delta = \delta^{-1}$ and eventually
$R_\delta = L_\delta+2$. On the same given slow $W$ define the full, unprojected
input
\begin{equation}\label{E:tube-actual-input}
  W_m(f)=W(T_{\delta,m}f),\qquad
  X_m(u)=g_\delta M_R(W_m)(\min(u,\ell_m)),\qquad
  g_\delta=\delta^{(1-a)/2}.
\end{equation}
Write $P_{\delta,E}^R$ for the physical $E$ flow driven by $g_\delta M_R(W\circ
T_{\delta,0})$ on the global fast horizon $[0,L_\delta]$. Let
$\mathscr{B}_\delta$ be the finite product of little
$C^\theta([0,B];\mathcal{K}(E))$ spaces of paths starting at zero, with exact
norm
\[
  \|x\|_{\delta,p}=\left(\delta\sum_{m=0}^N
  \|x_m\|_{C^\theta}^p\right)^{1/p}.
\]

\begin{proposition}[Eligibility of the E law]\label{P:tube-actual-law}
  The tuple $X_\delta = (X_m)$ satisfies the preceding Gaussian hypotheses with
  $H = H_I$ and
  \begin{align}
    (\mathcal{M}_\delta\Phi)_m(u)
    &=\int_0^{\min(u,\ell_m)}
    S_E(-v)\mathcal{T}_R\Phi(s_m+\delta v)S_E(v)\ud v,
    \label{E:tube-actual-mean}\\
    \E[X_\delta W(\Phi)]&=\sqrt\delta\,\mathcal{M}_\delta\Phi .
    \label{E:tube-actual-covariance}
  \end{align}
  Its support is $E_{0,\delta} = \overline{C_\delta(H_I)} \subset
  \mathscr{B}_\delta$, where $C_\delta\Phi = \E[X_\delta W(\Phi)] =
  \sqrt\delta\,\mathcal{M}_\delta\Phi$. The map $\mathcal{M}_\delta$ is compact
  at each fixed radius and parameters. At $R = R_\delta$,
  \begin{equation}\label{E:tube-typical-radius}
    b_\delta=2\E\|X_\delta\|_{\delta,p}
    \le C_{Q,B,p}\delta^{(1-a)/2-d/Q}\longrightarrow0 .
  \end{equation}
\end{proposition}

\begin{proof}
  Compact operators on separable Hilbert $E$ form a separable operator-norm
  space: finite-rank rational matrices in an orthonormal basis are dense.
  Polygonal paths give separability of the little H\"older closure. The
  construction is a $C^\theta$ limit of Lipschitz finite coefficient paths
  starting at zero; stopping at $\ell_m$ and extending constantly preserves this
  space and gives zero for an empty final interval. Gaussian transfer retains
  the given $W$ for all finitely many blocks. Every finite collection of linear
  statistics is an $L^2$ limit of joint first Wiener chaos Gaussians, hence
  jointly centered Gaussian. The finite tuple has finite $L^2$ norm by the
  fixed-parameter path moments in Corollary~\ref{C:radius-covariance-residuals}.
  The original Hilbert space $H_I$ is separable by
  Lemma~\ref{L:cameron-supported-tensor}. No separability of all
  $\mathcal{L}(E)$ or full $C^\theta$ is claimed.

  The $E$ mean field $\mathcal{T}_R$ is the compression of
  Proposition~\ref{P:cameron-full-mean}; it is compact into
  $L^{r_t}(I;\mathcal{K}(E))$ at fixed $R$.
  Proposition~\ref{P:uniform-mean-regulator} gives $\|\mathcal{T}_R\| \le C_{\rm
  mean} = M_0\|C_t\|$ independently of $R \ge 1$. Compact conjugation, affine
  time pullback and $1-1/r_t > \theta$ make the primitive
  in~\eqref{E:tube-actual-mean} bounded from that field space to
  $\mathscr{B}_\delta$, with constant stopping. Composition proves compactness
  of $\mathcal{M}_\delta$. There is no uniform-radius compactness assertion.

  The rank-zero identity~\eqref{E:cameron-actual-covariance}, compressed to $E$
  as in Proposition~\ref{P:symplectic-block-mean},
  gives~\eqref{E:tube-actual-covariance}. Its scale is exactly
  $g_\delta\delta^{a/2} = \sqrt\delta$. Every stopped row's slow pullback lies
  in $H_I$ by the inside-$I$ endpoint approximation of
  Lemma~\ref{L:cameron-actual-residual}. Countable matrix coordinates and
  rational times identify the tuple in the closed first Wiener chaos $W(H_I)$,
  so every linear statistic has its label there. Equivalently, apply that lemma
  at $J = 0$ on $F$ and compose any $E$-tuple statistic with the bounded
  $E$-compression. This is same-$W$ covariance, without a block-independence
  assumption or a time-indicator multiplier on the covariance space.
  Proposition~\ref{P:tube-wiener-realization} therefore applies to this exact
  law on its support. Finally, Proposition~\ref{P:radius-mean-anchor}, with zero
  residual projection, gives~\eqref{E:tube-typical-radius}. The third fixed-$Q$
  condition $Q > 2d/(1-a)$ makes this exponent positive; the first two supply
  the path moments and chaining. Fast directions used for the symplectic
  baseline do not replace this global $H_I$ action or get summed as separate
  costs. Zero covariance gives zero operator paths and means directly. This
  proves Proposition~\ref{P:tube-actual-law}.
\end{proof}

\begin{lemma}[The endpoint functional]\label{L:tube-actual-endpoint}
  For every $x \in \mathscr{B}_\delta$ define
  \[
    \begin{aligned}
      P_m(x_m)&=S_E(\ell_m)U_{x_m}(\ell_m,0),\qquad
      \mathcal{P}_\delta(x)=P_N(x_N)\cdots P_0(x_0),\\
      F_\delta(x)&=2\delta\log\max(1,\|\mathcal{P}_\delta(x)\|).
    \end{aligned}
  \]
  This is a finite continuous functional on all tuples. For the input it
  satisfies
  \begin{equation}\label{E:tube-global-endpoint}
    \mathcal{P}_\delta(X_\delta)
    =P_{\delta,E}^R(L_\delta,0),\qquad
    e^{F_\delta(X_\delta)/\delta}
    =\|P_{\delta,E}^R(L_\delta,0)\|^2 .
  \end{equation}
  These are almost-sure identities for each fixed admissible $\delta,R,B$,
  without a simultaneous event over all parameter choices.
\end{lemma}

\begin{proof}
  Generic Young evolution is globally defined and locally Lipschitz in every
  $C^\theta$ operator path. Finite products, fixed-time free factors, norm and
  $\log\max(1,\cdot)$ preserve continuity. Thus, $F_\delta$ is defined even off
  every Gaussian full-measure set. Every $x_m$ in this ambient space is a
  compact-valued $C^\theta$ path starting at zero. The finite-product assertion
  of Lemma~\ref{L:symplectic-free} gives, for every such tuple,
  \[
    \|\mathcal{P}_\delta(x)\|
    \ge\left\|S_E\Big(\sum_{m=0}^N\ell_m\Big)\right\|\ge1 .
  \]

  Lemma~\ref{L:raw-block-restart}, in its case without the deterministic mean,
  applies to this same full input with base physical profile $\eta^2$. To check
  the exact objects, use its single regularized slow noise preserving the given
  noise and pull it back by $T_{\delta,0}$ and $T_{\delta,m}$. At fixed $\delta$
  the global fast source random field is $\delta^{a/2}Y^\rho(\delta u,x)$ and
  its local restriction is its time shift by $mB$. Their ordinary Duhamel flows
  have the same massless free groups, amplitude $g_\delta$ and mask $m_R$.
  Source limits are taken on the fixed finite extended horizon $(N+1)B$ and on
  each local $B$ before composing the finitely many factors. The original
  cocycle from that lemma therefore gives the ordered physical $F$ product with
  the true remainder. The full operator paths and free groups intertwine through
  the closed isometry $i:E \to F$, including the source limits. Restricting the
  product and using injectivity of $i$ gives the first identity
  in~\eqref{E:tube-global-endpoint}. Both fast velocity and displacement remain
  in each factor. No mask was commuted through a free group, and the time blocks
  remain correlated on the same $W$. For zero covariance every input is zero and
  the ordered factors are $S_E(\ell_m)$, whose product is $S_E(L_\delta)$.

  The lower bound already established for every tuple removes the maximum at
  $X_\delta$. Together with the first identity, this proves the second identity
  in~\eqref{E:tube-global-endpoint}. No fixed-radius Gaussian-operator rate or
  source-to-point result is an input to this restart argument. This proves
  Lemma~\ref{L:tube-actual-endpoint}.
\end{proof}

\subsection{Uniform coarse envelope and the exact radial integral}

\begin{proposition}[Coarse tube envelope and operator moment]
  \label{P:tube-coarse-envelope}
  On the support $E_{0,\delta}$ let $A_\delta$ be its closed ball of radius
  $b_\delta$ and put
  \[
    G_\delta(r)=
    \sup_{\substack{a\in A_\delta\\\|\Phi\|_{H_I}\le r}}
    F_\delta(a+\mathcal{M}_\delta\Phi).
  \]
  This is finite, nondecreasing and Borel. There is $C_{\rm env} < \infty$,
  independent of $\delta$ and $R \ge 1$, such that
  \begin{equation}\label{E:tube-coarse-envelope}
    G_\delta(r)\le C_{\rm env}(1+b_\delta^\nu+r^\nu),\qquad
    \E\|P_{\delta,E}^R(L_\delta,0)\|^2
    \le\int_0^\infty\frac r\delta
    e^{(G_\delta(r)-r^2/2)/\delta}\ud r .
  \end{equation}
\end{proposition}

\begin{proof}
  Put $G_B = \sup_{|v|\le B}\|S_E(v)\|$ and $C_{\rm pr} =
  G_B^2[B^{1-1/r_t}+B^{1-1/r_t-\theta}]$. The lower bound $G_B \ge 1$ follows
  directly from $S_E(0) = I$. Proposition~\ref{P:uniform-actual-cameron} gives
  $1 \le G_B \le 3B$, $C_{\rm pr} \le 18B^{3-1/r_t}$. Its primitive estimate,
  now without the covariance factor $\sqrt\delta$, gives
  \begin{equation}\label{E:tube-primitive-sum}
    \delta\sum_m h_m^{r_t}
    \le(C_{\rm pr}C_{\rm mean})^{r_t}\|\Phi\|_{H_I}^{r_t},
    \qquad h_m=\|(\mathcal{M}_\delta\Phi)_m\|_{C^\theta}.
  \end{equation}
  Indeed, each primitive norm is at most $C_{\rm pr}\delta^{-1/r_t}
  \|\mathcal{T}_R\Phi\|_{L^{r_t}(I_m)}$. Raising to $r_t$ cancels the $\delta$
  weight, and the true partition sums to the full field norm, including an empty
  cell.

  The generic Banach Young bound supplies $c_Y = c_Y(B,\theta) > 0$ with
  $\log\|U_d(u,0)\| \le c_Y(1+\|d\|_{C^\theta}^{\nu})$ for $u \le B$. This
  follows from~\eqref{E:young-upper-refined-growth}
  and~\eqref{E:constant-young-growth} by $d \le 1+d^\nu$ and $(1+d)^\nu \le
  2^{\nu-1}(1+d^\nu)$. It applies on $\mathcal{L}(E)$ with multiplication norm
  one, without Gaussian, positivity or Hamiltonian assumptions on $d$. For $a
  \in A_\delta$ put $v_m = \|a_m\|_{C^\theta}$ and $n = N+1$. Convexity and
  submultiplicativity give
  \[
    F_\delta(a+\mathcal{M}_\delta\Phi)
    \le2\delta n(\log G_B+c_Y)
    +2^\nu c_Y\delta\sum_m(v_m^\nu+h_m^\nu).
  \]
  The upper expression is nonnegative, so it also bounds the logarithm with the
  maximum at one. The exact finite weighted power-mean inequalities are
  \begin{align*}
    \delta\sum_m v_m^\nu
    &\le(\delta n)^{1-\nu/p}\|a\|_{\delta,p}^\nu,\\
    \delta\sum_m h_m^\nu
    &\le(\delta n)^{1-\nu/r_t}
    \left(\delta\sum_m h_m^{r_t}\right)^{\nu/r_t}.
  \end{align*}
  They use exponents $p/\nu \ge 1$, $r_t/\nu > 1$ and nonnegative inputs
  $v_m^\nu,h_m^\nu$, followed by positive powers. Zero norms and empty cells
  cause no negative powers; $\delta > 0$ and $n \ge 2$, including the case $p =
  \nu$. Since $\delta n \le 2/B$, these bounds and~\eqref{E:tube-primitive-sum}
  give $F_\delta(a+\mathcal{M}_\delta\Phi) \le D_0+D_1b_\delta^\nu+D_2r^\nu$,
  where
  \begin{align*}
    D_0&=2(2/B)(\log G_B+c_Y),\\
    D_1&=2^\nu c_Y(2/B)^{1-\nu/p},\\
    D_2&=2^\nu c_Y(2/B)^{1-\nu/r_t}
    (C_{\rm pr}C_{\rm mean})^\nu .
  \end{align*}
  Hence, $C_{\rm env} = 1+D_0+D_1+D_2$ is independent of $\delta,R$ and proves
  the coarse bound for every admissible $\delta$, not just when $b_\delta$ is
  small. The supremum sets are nonempty and this bound makes their suprema
  finite. Monotonicity gives Borel measurability without asserting attainment of
  these suprema. The definition bounds $F_\delta$ on every closed tube. Apply
  Proposition~\ref{P:tube-rayleigh-transfer} and~\eqref{E:tube-global-endpoint}
  to obtain the operator-moment bound. For zero law its direct case applies,
  with no assumption that the free endpoint norm is one. This proves
  Proposition~\ref{P:tube-coarse-envelope}.
\end{proof}

\begin{lemma}[Exact Rayleigh tail bound]\label{L:tube-rayleigh-tail}
  Suppose $G$ is Borel and $G(r) \le A_0+Cr^\nu$, $C \ge 0$, $1 < \nu < 2$. If
  $M > 0$ and $4C \le M^{2-\nu}$, then
  \begin{equation}\label{E:tube-rayleigh-tail}
    \int_M^\infty\frac r\delta
    e^{(G(r)-r^2/2)/\delta}\ud r
    \le4e^{(A_0-M^2/4)/\delta}.
  \end{equation}
  In particular, the operator-moment bound in~\eqref{E:tube-coarse-envelope} is
  finite at fixed parameters.
\end{lemma}

\begin{proof}
  The finite scalar cutoff inequality gives, for $r \ge M$, $Cr^\nu-r^2/2 \le
  -M^2/8-r^2/8$. Its exact conditions are $M > 0$, $r \ge M$, $C \ge 0$, $1 <
  \nu < 2$ and $4C \le M^{2-\nu}$. Indeed, $r^{2-\nu} \ge M^{2-\nu} \ge 4C$
  gives $Cr^\nu \le r^2/4$, while $r^2 \ge M^2$ splits the remaining $-r^2/4$
  into the two displayed negative terms. Analytically, $-\frac{\ud}{\ud
  r}[4e^{-r^2/(8\delta)}] = (r/\delta)e^{-r^2/(8\delta)}$. Integrating gives the
  front factor $4e^{-M^2/(8\delta)}$; multiplication by $e^{(A_0-M^2/8)/\delta}$
  proves \eqref{E:tube-rayleigh-tail}. Such a cutoff exists since $2-\nu > 0$.
  On $[0,M]$ the growth bound gives $G(r) \le A_0+CM^\nu$, so the integral is
  finite there as well. Apply this to $A_0 = C_{\rm env}(1+b_\delta^\nu)$ and $C
  = C_{\rm env}$ in the preceding proposition. No infinite-dimensional Gaussian
  density is used. This proves Lemma~\ref{L:tube-rayleigh-tail}.
\end{proof}

\begin{lemma}[Supremum with the Gaussian quadratic cost]\label{L:tube-half-action}
  For finite $J:H \to \R$ bounded above on every original ball, let $g_J(r) =
  \sup_{\|h\|\le r}J(h)$. In the extended reals,
  \[
    \sup_{r\ge0}\{g_J(r)-r^2/2\}
    =\sup_{h\in H}\{J(h)-\|h\|_H^2/2\}.
  \]
\end{lemma}

\begin{proof}
  For the lower inequality take $r = \|h\|$. For the upper inequality, every
  $\|h\| \le r$ gives $J(h)-r^2/2 \le J(h)-\|h\|^2/2$. Take first the ball
  supremum and then the radius supremum. Each ball is nonempty with finite
  supremum, so no infinity-minus-infinity expression occurs. No compactness,
  rank or density assumption is needed. This proves
  Lemma~\ref{L:tube-half-action}.
\end{proof}

The Gaussian comparison retains one half of the squared norm of the single
control. The coarse envelope is not the sharp spectral/product envelope. Its
application to the original canonical point moment, and the sharp envelope and
variational comparison, require their separate proofs.
